\begin{prob}
    Provide a tight theoretical lower bound for the problems given below. Provide justification for your answer.

    \begin{subprobset}

        \begin{subprob}
            Given a list of size $n$ containing \python{True}s and \python{False}s, determine whether \python{True} or
            \python{False} is more common (or if there is a tie).

            \begin{soln}
                $\Omega(n)$.

                In the worst case, we need to read the whole list, as the ``winner'' will be determined by the very last entry.

                Theoretical lower bound concerns the worst case, but note that even the best case will require us to read at least
                half of the entries here. For example, if $n = 100$ and the first 51 entries are \python{True}, we can conclude that
                \python{True} is the winner after reading those 51. But if even one of them is \python{False}, it may be that
                \python{False} is the winner if it makes up all of the remaining entries.
            \end{soln}

        \end{subprob}

        \begin{subprob}
            Given a list of $n$ numbers, all assumed to be integers between 1 and 100, sort them.

            \begin{soln}
                $\Omega(n)$.

                At the minimum, we have to read all of the numbers, which takes $\Theta(n)$ time.

                You may have heard that ``sorting takes $\Theta(n \log n)$ time'', but we have to be careful. First,
                this applies only to \emph{comparison sorts}, where we sort by comparing elements of the input to one another.
                If we don't have any assumptions on the numbers, we have to do a comparison sort. But since we've assumed that
                they are between 0 and 100, we can actually sort them without ever comparing any two numbers. Here's the algorithm:

                \begin{enumerate}
                    \item Initialize a list, \python{counts}, with 100 entries, all zero to start.
                    \item Loop through the input list. For each number, $x$, increment $\python{counts[x]}$ by one. For example, if we see
                        7, increment \python{counts[7]} by one. In short, this loop counts how many times we see each number.
                    \item Finally, loop through the range $0, 1, \ldots, 100$, and print the number $x$ a number of times equal to \python{counts[x]}.
                        That is, if \python{counts[7]} is 3, print 3 sevens.
                \end{enumerate}

                Because we're looping through the input list once, and it has $n$ elements, we've sorted the list in $\Theta(n)$ time.
             \end{soln}

        \end{subprob}

        \begin{subprob}
            Given an $\sqrt{n} \times n$ array whose rows are sorted (but whose columns may not be), find the largest overall entry in the array.

            For example, the array could look like:

            \[
            \begin{pmatrix}
                -2 & 4 & 7 & 8 & 10 & 12 & 20 & 21 & 50\\
                -30 & -20 & -10 & 0 & 1 & 2 & 3 & 21 & 23\\
                -10 & -2 & 0 & 2 & 4 & 6 & 30 & 31 & 35
            \end{pmatrix}
            \]

            This is an $\sqrt n \times n$ array, with $n = 9$ (there are 3 rows and 9 columns). Each row is sorted, but the columns aren't.

            \begin{soln}
                $\Omega(\sqrt n)$

                The largest element of the array has to be in the last column, so we simply look through all of the entries in that column.
                There are $\sqrt n$ such entries, so we take $\Theta(\sqrt n)$ time.

                Note that we don't have to read all of the entries of the array -- we can actually ignore almost all of them.
            \end{soln}

        \end{subprob}

    \end{subprobset}

\end{prob}
